\chapter{Parcours client}

\section{Tableau de bord}

Le tableau de bord client synthétise les sessions actives et terminées, l'énergie
reçue, les montants réglés, les abonnements et les identifiants de recharge. Il
permet d'accéder rapidement à la recherche d'une station ou à une session active.

\section{Trouver une station}

La page \textbf{Find a station} propose une recherche par nom, ville ou adresse,
ainsi que plusieurs représentations. Les vues carte, cartes et liste servent des
besoins différents :

\begin{itemize}[leftmargin=1.6em]
  \item la carte localise les stations et leur état ;
  \item les cartes facilitent le choix sur mobile ;
  \item la liste compare rapidement disponibilité, puissance et connecteurs.
\end{itemize}

\begin{procedurebox}{Rechercher à proximité}
  \begin{enumerate}[leftmargin=1.6em]
    \item ouvrir \textbf{Find a station} ;
    \item sélectionner \textbf{Near me} ;
    \item autoriser la géolocalisation du navigateur ;
    \item appliquer le rayon défini dans \textbf{Settings} ;
    \item filtrer si nécessaire par disponibilité ou type de connecteur ;
    \item ouvrir la station retenue.
  \end{enumerate}
\end{procedurebox}

Depuis une station, le client peut copier les coordonnées, ouvrir l'emplacement
dans Google Maps et demander un itinéraire. Le bouton \textbf{Charge} ouvre
directement l'étape de sélection du connecteur lorsque la station est déjà
connue.

\section{Scanner un QR code}

Le QR code physique placé sur un connecteur contient une référence reconnue par
ChargeTrackr. L'action \textbf{Scan QR} du client sert à éviter une saisie
manuelle sur le terrain.

\begin{procedurebox}{Démarrer depuis le QR code physique}
  \begin{enumerate}[leftmargin=1.6em]
    \item ouvrir \textbf{Scan QR} depuis la recherche ;
    \item autoriser l'accès à la caméra ;
    \item cadrer le QR code du connecteur physique ;
    \item vérifier la station, le libellé et le type de prise proposés ;
    \item poursuivre le parcours de recharge.
  \end{enumerate}
\end{procedurebox}

\attention{Un QR code affiché dans les vues de gestion sert à imprimer ou
remplacer l'étiquette physique. Le client doit scanner l'étiquette présente sur
la borne et non une capture non vérifiée.}

\section{Abonnements et réseaux}

Un client peut utiliser plusieurs organisations. Un paiement ponctuel de recharge
ne l'inscrit pas automatiquement à l'organisation propriétaire de la borne.
L'adhésion à un plan est un acte distinct, volontaire et révocable.

\manualscreenshot{client-subscriptions.png}{Plans, réseaux et abonnements du client}{fig:user-client-subscriptions}

\begin{procedurebox}{Souscrire à un plan}
  \begin{enumerate}[leftmargin=1.6em]
    \item ouvrir \textbf{Plans \& subscriptions} ;
    \item comparer les réseaux, tarifs, réductions et conditions ;
    \item ouvrir le plan souhaité ;
    \item vérifier l'organisation et son logo ;
    \item confirmer le paiement simulé de l'abonnement ;
    \item contrôler l'apparition du plan dans \textbf{My memberships}.
  \end{enumerate}
\end{procedurebox}

\expected{Le plan devient actif pour l'organisation choisie. Ses avantages sont
appliqués aux futures recharges éligibles, sans modifier les autres adhésions.}

\section{Démarrer une recharge}

Le parcours est guidé en cinq étapes. Lorsqu'une station a déjà été choisie,
l'application passe directement à l'étape \textbf{Connector}.

\begin{longtable}{|C{1cm}|L{3.1cm}|L{5.2cm}|L{5.8cm}|}
  \hline
  \rowcolor{ctmint}\textbf{\#} & \textbf{Étape} & \textbf{Dans l'application} &
  \textbf{Dans le monde physique} \\ \hline
  \endfirsthead
  \hline
  \rowcolor{ctmint}\textbf{\#} & \textbf{Étape} & \textbf{Dans l'application} &
  \textbf{Dans le monde physique} \\ \hline
  \endhead
  1 & Station & Rechercher ou confirmer la station. & Se placer devant la borne
  identifiée et vérifier sa référence. \\ \hline
  2 & Connector & Choisir un connecteur disponible et compatible. & Repérer le
  libellé physique, par exemple A1, et le type de prise. \\ \hline
  3 & Connect & Attendre la détection OCPP automatique. & Stationner, couper le
  véhicule et brancher complètement le câble. \\ \hline
  4 & Payment & Autoriser le moyen de paiement et consulter le tarif. & Aucune
  action sur la borne ; l'autorisation est traitée dans l'application. \\ \hline
  5 & Start & Confirmer le démarrage distant. & Attendre le verrouillage du câble
  et l'indication de charge sur le véhicule ou la borne. \\ \hline
\end{longtable}

\expected{Le connecteur passe à l'état occupé, une session active est créée et les
mesures reçues par OCPP actualisent durée, énergie, puissance et estimation.}

\section{Paiement de la recharge}

Dans le MVP du stage, le paiement utilise un adaptateur externe simulé. Le
parcours reproduit les étapes attendues d'un prestataire réel :

\begin{enumerate}[leftmargin=1.6em]
  \item autorisation avant le démarrage ;
  \item référence fournisseur et réponse de webhook ;
  \item calcul du montant avec le tarif et l'abonnement ;
  \item règlement final à l'arrêt ;
  \item génération d'un reçu consultable et téléchargeable.
\end{enumerate}

Le simulateur n'effectue aucun débit bancaire réel. L'architecture permet de
remplacer ultérieurement l'adaptateur par un fournisseur tel que ClicToPay,
e-DINAR/D17, GPGCheckout, Konnect ou Flouci.

\section{Suivre et arrêter la session}

Une vue en direct présente les mesures et propose :

\begin{itemize}[leftmargin=1.6em]
  \item \textbf{Continue in background} pour fermer la fenêtre tout en gardant la
  session active ;
  \item \textbf{Open live view} pour rouvrir la vue ;
  \item \textbf{Stop charging} pour demander l'arrêt.
\end{itemize}

Une session peut également s'arrêter lorsque le véhicule ou la borne atteint sa
condition de fin, en cas d'erreur de sécurité, de perte de communication
confirmée ou d'arrêt distant autorisé par l'exploitation.

\begin{procedurebox}{Arrêter volontairement une recharge}
  \begin{enumerate}[leftmargin=1.6em]
    \item ouvrir \textbf{My sessions} ;
    \item ouvrir la vue de la session active ;
    \item sélectionner \textbf{Stop charging} ;
    \item confirmer la demande ;
    \item attendre l'état terminé avant de débrancher ;
    \item consulter le montant final et le reçu.
  \end{enumerate}
\end{procedurebox}

\attention{Ne jamais forcer le retrait d'un câble verrouillé. Attendre la
confirmation d'arrêt de la borne et du véhicule.}

\manualscreenshot{client-sessions.png}{Suivi des sessions de recharge du client}{fig:user-client-sessions}
